\section{Relative Perron envelopes of arbitrary strong completed families}
\label{sec:r11-strong-completed-Perron}

Retain exactly the strict analytic box, background extension
and full strong first-energy class of
Section~\ref{sec:r11-strong-completed-comparison}.
The background is the unchanged arbitrary original symmetric
middle-positive separately closed $S$.
Write $f\preceq g$ if $f\leq g$ locally outside an
ordinary-open outer $S$-capacity-null set. This order concerns
the completed Borel representatives. No ordinary ambient
upper regularization is used.

The frozen R2 relative capacitary envelope theorem is
\texttt{INTERNAL\_INPUT} in its bounded scope. Theorem
\ref{thm:r11-barrier-relative-Perron} extends an actual family
above one prescribed actual power-energy barrier. The theorem
below instead permits an arbitrary nonempty family of strong
completed first-energy potentials, with no common actual lower
barrier. Its countable selection is proved directly by global
energy minimization over finite maxima.

\begin{theorem}[The full strong completed relative envelope]
\label{thm:r11-strong-completed-Perron}
Let $\mathscr F\subset\widehat{\mathcal E}^1_{\rm str}(V)$
be any nonempty family, possibly uncountable. Let
$\mathscr F^\vee$ consist of all nonempty finite maxima of
its members, and put
\[
                    L=\inf_{w\in\mathscr F^\vee}\widehat e(w).
 \tag{SE1}
\]
There is a unique relative class
$H\in\widehat{\mathcal E}^1_{\rm str}(V)$ with the following
least-upper-bound properties:
\[
 \begin{split}
 &f\preceq H\qquad(f\in\mathscr F),\\
 &\bigl[f\preceq g\text{ for every }f\in\mathscr F,
                      \ g\text{ Borel}\bigr]
                                      \Longrightarrow H\preceq g.
 \end{split}
 \tag{SE2}
\]
The exceptional set in the first line may depend on $f$;
there is no uncountable-union nullity assertion.

There are countably many selected members of $\mathscr F$
whose cumulative finite maxima $v_j$ satisfy
\[
                  v_j\uparrow H,\qquad
                  \widehat I(v_j,H)\longrightarrow0,
                  \widehat e(H)=L.
 \tag{SE3}
\]
Here the actual increasing Borel limit may be used as the
representative of $H$. More generally, every sequence
$w_j\in\mathscr F^\vee$ with $\widehat e(w_j)\to L$
converges strongly to $H$ and therefore has all the critical
current and mixed-energy limits (IC11)--(IC13).
\end{theorem}

\noindent\textbf{Status: \texttt{NEW\_PROVED}} for arbitrary
completed families and their full energy boundary class.

\begin{proof}
Finite maxima are members of the strong completion by (CC1).
Their energies are nonnegative, and one fixed family member
has finite energy. Thus $0\leq L<\infty$.
Choose $w_j\in\mathscr F^\vee$ with
$\widehat e(w_j)\leq L+1/j$, and put
\[
                            v_j=\max(w_1,\ldots,w_j).
 \tag{SE4}
\]
Each is still a finite maximum of members of the given family.
It increases with $j$, and energy monotonicity gives
$L\leq\widehat e(v_j)\leq\widehat e(w_j)$.
The increasing-closure theorem
\ref{cor:r11-strong-monotone-closure} therefore gives its
actual Borel limit $H$ in the strong completion, with strong
convergence and $\widehat e(H)=L$.
Only countably many members of the given family occur in all the
selected $w_j$; this proves the countable assertion in (SE3).

Fix an arbitrary $f\in\mathscr F$ and set $z_j=\max(v_j,f)$.
Then $z_j\in\mathscr F^\vee$ and
\[
                       L\leq\widehat e(z_j)
                                      \leq\widehat e(v_j).
\]
The strong maximum continuity (CC2) gives
$z_j\to z:=\max(H,f)$ strongly. Hence
$\widehat e(z)=L=\widehat e(H)$.
Since $H\leq z$, the ordered discrepancy estimate (CI1)
gives
\[
               0\leq\widehat I(H,z)
                      \leq\widehat e(H)-\widehat e(z)=0.
 \tag{SE5}
\]
Separation of the completed classes implies $z=H$
locally outside an outer-null set. Thus $f\preceq H$.
This argument fixes one $f$ at a time, as required by (SE2).

Now let $g$ be any Borel upper bound in the sense of the
second line of (SE2). It bounds each of the countably many
selected family members outside its corresponding null set.
Their countable union is locally outer-null. Off that union,
$g$ bounds every $v_j$ and hence their actual pointwise
limit $H$. This proves leastness among Borel upper bounds,
without requiring $g$ to be psh or completed.
Two such least upper bounds bound each other, so their
relative classes agree.

Finally, every $w\in\mathscr F^\vee$ satisfies $w\preceq H$,
by the first line of (SE2) and finite subadditivity.
For any sequence whose energies tend to $L$, (CI1) gives
\[
             0\leq\widehat I(w_j,H)
                    \leq\widehat e(w_j)-\widehat e(H)
                                      \longrightarrow0.
 \tag{SE6}
\]
The strong continuity theorem supplies all the current and
critical weighted-measure limits claimed in the statement.
\end{proof}

\begin{corollary}[Subsolution envelopes and Borel upper constraints]
\label{cor:r11-strong-Perron-subsolution}
Let $\nu$ be a fixed positive Radon measure on $V$ and
$\psi$ a Borel upper constraint. If the preceding nonempty
family satisfies
\[
                    \widehat\mu_f\geq\nu,
                    \qquad f\preceq\psi
                                  \quad(f\in\mathscr F),
 \tag{SE7}
\]
then its canonical relative envelope satisfies
\[
                             \widehat\mu_H\geq\nu,
                             \qquad H\preceq\psi.
 \tag{SE8}
\]
The current inequality holds on every original compact smooth
scalar test and every actual moment refinement. A finite total
mass for $\nu$ is not needed; the printed Radon hypothesis is
local. No equality $\widehat\mu_H=\nu$ or contact support
assertion is inferred.
\end{corollary}

\noindent\textbf{Status: \texttt{NEW\_PROVED}} for full strong
completed subsolution families; it is a consequence of the
new finite-maximum and envelope theorems.

\begin{proof}
The completed weak maximum inequality, proved with its full
middle-positive cone and equality-side mass, shows that every
finite maximum of members of $\mathscr F$ has top measure at
least $\nu$. In particular $\widehat\mu_{v_j}\geq\nu$ for
the selected increasing sequence (SE4).
Its strong convergence passes every original compact scalar
test, so positivity gives $\widehat\mu_H\geq\nu$.
All measures are locally Radon; no global mass subtraction
or unproved boundary tightness is needed for this inequality.
The least-upper-bound property (SE2), or the countable selected
members directly, gives $H\preceq\psi$.
\end{proof}

\paragraph{The distinction from an ambient Perron construction.}
This theorem supplies a canonical relative least upper bound
and a retained source lower bound inside the complete strong
class. The ordinary upper semicontinuous envelope can still
lose the relative mass, as in the all-power counterexample
(PO1)--(PO7). Moreover the exact strong completed compact-source
obstruction (SO1)--(SO6) prevents promoting (SE8) to a universal
source equation. These are separate issues from the proved
existence of this Borel relative envelope. General raw-capacity
outer regularity, unrestricted ordinary-boundary Dirichlet
existence, GQ and low-real-degree CTO remain outside the claim.
